\chapter{Introdução}

A AstraZeneca PLC é uma empresa farmacêutica global, orientada pela ciência, cotada em Londres, Estocolmo e Nova Iorque \cite[p.~71]{AR}, com \$58,739m de Receita Total \cite[p.~125]{AR} e 96,100 empregados em 2025 \cite[p.~27]{AR}. Esta secção apresenta o que a empresa é, o que faz, onde opera e qual a sua dimensão, e termina com o seu contributo para os Objetivos de Desenvolvimento Sustentável da ONU.

{\small\color{azmuted}
Salvo indicação em contrário, os dados provêm do \textit{Annual Report and Form 20-F Information 2025} da AstraZeneca (``o Relatório e Contas'', abreviado ``AR'' no texto) \cite{AR}; os números de página são os que estão impressos nas páginas do relatório. A AstraZeneca PLC e as suas subsidiárias são designadas por ``o Grupo'', os valores estão em dólares norte-americanos, e ``\$m'' e ``\$bn'' significam milhões e milhares de milhões. Os blocos marcados como Análise são interpretação dos autores, não afirmações do Relatório e Contas.
\par}

\section{Tipo de organização}

A AstraZeneca PLC é uma \textit{public limited company}, de responsabilidade limitada por ações, constituída e domiciliada em Inglaterra e no País de Gales \cite[p.~199]{AR}. Foi constituída em 17 de junho de 1992 ao abrigo do Companies Act 1985, e a sua sede social e sede corporativa situam-se em Cambridge, Reino Unido \cite[Item~4.A]{F20}, \cite[p.~230]{AR}. Em termos portugueses corresponde a uma \textit{sociedade anónima} cotada, uma forma jurídica que o material da unidade curricular cita entre as empresas privadas \cite[p.~5]{INT}.

\begin{itemize}[leftmargin=*]
\item \textbf{1992--1993.} Empresa constituída em 1992. Em 1993 foram cindidos os negócios farmacêutico, agroquímico e de químicos de especialidade da Imperial Chemical Industries PLC; a empresa chamou-se Zeneca Group PLC até abril de 1999. \cite[Item~4.A]{F20}
\item \textbf{1999.} Fusão com a Astra da Suécia; a empresa mudou o nome para AstraZeneca PLC. \cite[Item~4.A]{F20}
\item \textbf{2000.} O negócio agroquímico foi cindido e fundido com um negócio semelhante da Novartis para formar a Syngenta AG. \cite[Item~4.A]{F20}
\item \textbf{2007.} Aquisição da MedImmune, um negócio de biológicos e vacinas sediado nos Estados Unidos. \cite[Item~4.A]{F20}
\item \textbf{2021.} Aquisição da Alexion, um negócio de doenças raras sediado nos Estados Unidos. \cite[Item~4.A]{F20}
\end{itemize}

O material da unidade curricular classifica as organizações segundo quem beneficia delas. A AstraZeneca insere-se no tipo de \textbf{interesse comercial}, que é habitualmente uma empresa privada cujo benefício principal reverte para os seus proprietários \cite[pp.~4--5]{INT}, e no tipo \textbf{privado} quando classificada segundo quem detém o capital \cite[p.~6]{INT}. É uma empresa \textbf{com fins lucrativos}: em 2025 obteve um resultado do período de \$10,233m e pagou dividendos de \$4,846m \cite[p.~125]{AR}; para 2025, o Conselho declarou um dividendo anual de \$3.20 por ação \cite[p.~3]{AR}, estimado em \$4,963m no total \cite[p.~225]{AR}.

As suas ações estão cotadas na London Stock Exchange, na Nasdaq Stockholm e na New York Stock Exchange \cite[p.~71]{AR}. Desde 2 de fevereiro de 2026 as ações ordinárias estão cotadas diretamente na NYSE, substituindo a anterior cotação nos Estados Unidos de American Depositary Shares na Nasdaq \cite[p.~50]{AR}.

\section{Detentores do capital}

A 31 de dezembro de 2025 o capital social era composto por 1,550,907,927 ações ordinárias de \$0.25 cada, mais 50,000 ações preferenciais resgatáveis \cite[p.~224]{AR}. O AR afirma que os principais acionistas não têm direitos de voto diferentes \cite[p.~225]{AR} e descreve uma base acionista geograficamente diversificada \cite[p.~76]{AR}. Quatro grupos detinham 3\% ou mais.

\begin{xltabular}{\linewidth}{@{}Yrrr@{}}
\caption{Acionistas com 3\% ou mais dos direitos de voto. A Wellington Management Company LLP, uma subsidiária, está incluída no valor do Wellington Management Group. Fonte: \cite[p.~225]{AR}.}\label{tab:acionistas}\\*
\toprule
\textbf{Acionista} & \textbf{Ações na última comunicação} & \textbf{\% na última comunicação} & \textbf{\% a 31 dez. 2025} \\*
\midrule
\endfirsthead
\toprule
\textbf{Acionista} & \textbf{Ações na última comunicação} & \textbf{\% na última comunicação} & \textbf{\% a 31 dez. 2025} \\*
\midrule
\endhead
\bottomrule
\endfoot
BlackRock & 100,885,181 & 6.96 & \textbf{6.50} \\
The Capital Group Companies & 77,125,348 & 4.97 & \textbf{4.97} \\
Wellington Management Group LLP & 65,120,892 & 4.96 & \textbf{4.20} \\
Investor AB & 51,587,810 & 3.93 & \textbf{3.33} \\
\end{xltabular}

A maior participação a 31 de dezembro de 2025 é portanto de 6.50\%, e a Empresa não foi notificada de qualquer outra pessoa que detenha uma participação sujeita a comunicação.

\section{Atividades e âmbito}

A AstraZeneca descreve-se como um negócio farmacêutico global, orientado pela ciência e centrado no doente, comprometido com a excelência na investigação, desenvolvimento, fabrico e comercialização de medicamentos sujeitos a receita médica nas áreas de cuidados primários, cuidados especializados e doenças raras \cite[p.~8]{AR}.

A organização comercial está dividida em três regiões: os Estados Unidos, a Europa-Canadá e a Internacional, que inclui a Austrália, a Nova Zelândia e os Mercados Emergentes (entre eles a China). O Japão reporta separadamente \cite[p.~32]{AR}. Os produtos chegam ao mercado através de empresas locais de marketing detidas na totalidade, de distribuidores e de escritórios de representação locais \cite[p.~32]{AR}.

\begin{itemize}[leftmargin=*]
\item \textbf{Investigação e desenvolvimento --- 16,100:} empregados em I\&D nos locais globais do Grupo \cite[p.~27]{AR}
\item \textbf{Comercial --- 47,400:} empregados na área comercial; a rede tem 10 \textit{hubs} globais \cite[p.~27]{AR}
\item \textbf{Fabrico --- 16,900:} empregados no fabrico, em 31 locais de operações em 15 países \cite[p.~27]{AR}
\item \textbf{Total --- 96,100:} empregados (2024: 94,300) \cite[p.~27]{AR}
\end{itemize}

\section{Principais produtos}

A receita está organizada em três áreas terapêuticas. A Oncologia é a maior, seguida dos BioFarmacêuticos, que agrupam a área cardiovascular, renal e do metabolismo, a respiratória e a imunologia, e as vacinas e terapias imunitárias \cite[pp.~2, 4]{AR}. As Doenças Raras surgem em terceiro lugar.

\begin{table}[H]\centering
\caption{Receita Total por área terapêutica em 2025. As percentagens somam 101\% por efeito de arredondamento. Fonte: \cite[p.~2]{AR}.}
\label{tab:receita-area}
\begin{tabular}{@{}lrrB@{}}
\toprule
\textbf{Área terapêutica} & \textbf{\$bn} & \textbf{\%} & \\
\midrule
Oncologia           & 25.6 & 44 & \barra{1.000} \\
BioFarmacêuticos    & 23.0 & 39 & \barra{0.898} \\
Doenças Raras       &  9.1 & 16 & \barra{0.355} \\
Outros Medicamentos &  1.0 &  2 & \barra{0.039} \\
\midrule
\textbf{Total}      & \textbf{58.7} & \textbf{100} & \\
\bottomrule
\end{tabular}
\end{table}
O AR indica as vendas das marcas mais vendidas. Em conjunto, os cinco maiores produtos geraram \$30,045m, o que corresponde a 51.1\% da Receita Total (calculado a partir de \cite[p.~56]{AR}: 30,045 ÷ 58,739). O Grupo conta 16 produtos com vendas anuais acima de \$1bn \cite[p.~50]{AR}.

\begin{table}[H]\centering
\caption{Cinco produtos mais vendidos em 2025. Fonte: \cite[p.~56]{AR}.}
\label{tab:top5}
\begin{tabular}{@{}lrB@{}}
\toprule
\textbf{Produto} & \textbf{\$m} & \\
\midrule
Farxiga   & 8,492 & \barra{1.000} \\
Tagrisso  & 7,254 & \barra{0.854} \\
Imfinzi   & 6,063 & \barra{0.714} \\
Ultomiris & 4,718 & \barra{0.556} \\
Calquence & 3,518 & \barra{0.414} \\
\bottomrule
\end{tabular}
\end{table}

\begin{table}[H]\centering
\caption{Marcas mais vendidas por área terapêutica. O AR associa a queda do Soliris à transição de doentes para o Ultomiris \cite[pp.~24, 56]{AR}.}
\label{tab:marcas}
\begin{tabular}{@{}llrr@{}}
\toprule
\textbf{Área} & \textbf{Produto} & \textbf{\$m} & \textbf{Variação} \\
\midrule
Oncologia \cite[p.~13]{AR}      & Tagrisso  & 7,254 & +10\% \\
                                & Imfinzi   & 6,063 & +29\% \\
                                & Calquence & 3,518 & +12\% \\
                                & Lynparza  & 3,279 & $-$11\% \\
\addlinespace
Doenças Raras \cite[p.~23]{AR}  & Ultomiris & 4,718 & +20\% \\
                                & Soliris   & 1,837 & $-$29\% \\
                                & Strensiq  & 1,678 & +19\% \\
                                & Koselugo  &   662 &  +5\% \\
\bottomrule
\end{tabular}
\end{table}
\section{Principais mercados}

Os Estados Unidos e a Europa representam em conjunto 65\% da Receita Total \cite[p.~2]{AR}. Os Mercados Emergentes são a segunda maior região e a que cresce mais depressa. ``Established RoW'' abrange o Japão, o Canadá, a Austrália e a Nova Zelândia \cite[p.~33]{AR}.

\begin{table}[H]\centering
\caption{Receita Total por região em 2025, com o crescimento a taxas de câmbio efetivas. As quotas são de 43\%, 26\%, 22\% e 9\% \cite[p.~2]{AR}; os valores provêm de \cite[p.~32]{AR} e somam \$58,739m.}
\label{tab:receita-regiao}
\begin{tabular}{@{}lrrB@{}}
\toprule
\textbf{Região} & \textbf{\$m} & \textbf{Crescimento} & \\
\midrule
Estados Unidos      & 25,450 & +10\% & \barra{1.000} \\
Mercados Emergentes & 15,303 & +12\% & \barrab{0.601} \\
Europa              & 12,739 &  +5\% & \barra{0.501} \\
Established RoW     &  5,247 &  +5\% & \barrab{0.206} \\
\bottomrule
\end{tabular}
\end{table}
\section{Países de presença}

A AstraZeneca tem presença ativa em mais de 80 países e vende os seus produtos em mais de 125 \cite[p.~32]{AR}. A sua rede inclui seis centros estratégicos de I\&D (Boston, Gaithersburg, Cambridge, Gotemburgo, Xangai e Pequim, este último inaugurado em 2025), dez \textit{hubs} globais e 31 locais de operações em 15 países \cite[p.~27]{AR}.

\figura{fig04-alcance-global}{1.0}{Alcance e presença global: centros de I\&D, \textit{hubs} globais, locais de operações e empregados por região de reporte. Fonte: \cite[p.~27]{AR}.}

\begin{table}[H]\centering
\caption{Empregados por região de reporte no final de 2025, num total de 96,100. Fonte: \cite[p.~27]{AR}.}
\label{tab:empregados}
\begin{tabular}{@{}lrrB@{}}
\toprule
\textbf{Região} & \textbf{Empregados} & \textbf{\%} & \\
\midrule
Mercados Emergentes & 36,800 & 38 & \barra{1.000} \\
Europa              & 34,000 & 35 & \barrab{0.924} \\
Estados Unidos      & 17,900 & 19 & \barra{0.486} \\
Established RoW     &  7,400 &  8 & \barrab{0.201} \\
\bottomrule
\end{tabular}
\end{table}
\section{Indicadores principais}

\begin{table}[H]\centering
\caption{Indicadores principais de 2025. Fontes: \cite[p.~1]{AR}; \cite[p.~125]{AR}; \cite[p.~27]{AR}.}
\label{tab:indicadores}
\begin{tabular}{@{}lr@{}}
\toprule
\textbf{Indicador (2025)} & \textbf{Valor} \\
\midrule
Receita Total                      & \$58,739m (+9\%) \\
\quad Vendas de Produtos           & \$55,573m \\
\quad Receita de Alianças          & \$3,067m \\
\quad Receita de Colaborações      & \$99m \\
\addlinespace
Despesa em I\&D                    & \$14,232m \\
Resultado operacional (reportado)  & \$13,743m \\
Resultado do período               & \$10,233m \\
Dividendos declarados e pagos      & \$4,846m \\
Empregados                         & 96,100 \\
\bottomrule
\end{tabular}
\end{table}

\begin{table}[H]\centering
\caption{Receita Total de 2023 a 2025, em \$bn. Fonte: \cite[p.~11]{AR}.}
\label{tab:receita-3anos}
\begin{tabular}{@{}lrB@{}}
\toprule
\textbf{Ano} & \textbf{\$bn} & \\
\midrule
2023 & 45.8 & \barra{0.780} \\
2024 & 54.1 & \barra{0.920} \\
2025 & 58.7 & \barra{1.000} \\
\bottomrule
\end{tabular}
\end{table}
\section{Contributo para os Objetivos de Desenvolvimento Sustentável da ONU}

O AR menciona os ODS da ONU uma única vez, quando afirma que as parcerias com organizações multilaterais e ONG os apoiam \cite[p.~76]{AR}. Não faz o mapeamento das ações para objetivos individuais. O mapeamento que segue é, por isso, Análise dos autores: cada objetivo é ligado às ações e metas que o Grupo reporta.

\subsection{O que o Grupo já faz}

\begin{xltabular}{\linewidth}{@{}c>{\hsize=1.45\hsize}Y>{\hsize=0.55\hsize}Y@{}}
\caption{ODS já apoiados pelo Grupo.}\label{tab:ods-atuais}\\*
\toprule
\textbf{ODS} & \textbf{Ação e evidência} & \textbf{Porque contribui} \\*
\midrule
\endfirsthead
\toprule
\textbf{ODS} & \textbf{Ação e evidência} & \textbf{Porque contribui} \\*
\midrule
\endhead
\bottomrule
\endfoot
3 & \textbf{Saúde de qualidade.} 320 milhões de pessoas impactadas positivamente desde 2024, 156 milhões delas de grupos carenciados. Ambição para 2030: mil milhões de pessoas, 400 milhões de comunidades carenciadas. Uma política global de preços escalonados com quatro escalões baseados no rendimento nacional bruto, e Programas de Assistência ao Doente. \cite[pp.~5, 41]{AR} & O negócio central é descobrir e fornecer medicamentos, e estes programas ampliam o acesso ao tratamento. \\
13 & \textbf{Ação climática.} As emissões de âmbito 1 e 2 caíram 88.1\% desde 2015, face a uma meta de 98\% até 2026. Metas: $-$50\% do âmbito 3 até 2030 e $-$90\% até 2045 (base de 2019), e zero emissões líquidas até 2045. \cite[pp.~2, 42]{AR} & Define metas de redução mensuráveis, verificadas no âmbito da Science Based Targets initiative \cite[p.~42]{AR}. \\
6 & \textbf{Água potável e saneamento.} O sítio do Grupo reporta uma redução de 23\% na captação de água face a 2015 \cite{NATURE}. O AR indica uma redução de 24\% no consumo de água apenas para a função de Operações \cite[p.~33]{AR}. & Reduz o uso de um recurso partilhado nas suas próprias operações. \\
12 & \textbf{Produção e consumo sustentáveis.} Os resíduos caíram 23\% desde 2015, com o objetivo de uma taxa de circularidade de resíduos de 50\% até 2030 \cite{NATURE}. O AR também indica $-$23\% de resíduos para as Operações \cite[p.~33]{AR}. & Corta resíduos na origem e define um objetivo de circularidade. \\
15 & \textbf{Proteger a vida terrestre.} AZ Forest, um programa de reflorestação e biodiversidade: 66 milhões de árvores plantadas em seis continentes desde 2020 (no final de 2025). O sítio indica uma meta de mais de 100 milhões até 2030 e que a meta foi revista em 2026 \cite{FOREST}. & Restaura ecossistemas e biodiversidade. \\
\end{xltabular}

\subsection{Objetivos adicionais sustentados por evidência do AR}

\begin{xltabular}{\linewidth}{@{}cY@{}}
\caption{Outros ODS sustentados por evidência no Relatório e Contas.}\label{tab:ods-adicionais}\\*
\toprule
\textbf{ODS} & \textbf{Evidência} \\*
\midrule
\endfirsthead
\toprule
\textbf{ODS} & \textbf{Evidência} \\*
\midrule
\endhead
\bottomrule
\endfoot
5 & \textbf{Igualdade de género.} As mulheres são 54\% da força de trabalho, há sete mulheres no Conselho de Administração (50\%), e cinco dos dez membros da Senior Executive Team são mulheres \cite[p.~39]{AR}. \\
17 & \textbf{Parcerias para a implementação dos objetivos.} Mais de 1,000 colaborações em todo o mundo \cite[p.~37]{AR}, e parcerias com organizações multilaterais e ONG em apoio dos ODS \cite[p.~76]{AR}. \\
\end{xltabular}

\subsection{O que o Grupo poderia fazer}

O próprio AR mostra onde estão as maiores lacunas. As três propostas seguintes partem dessas lacunas e não são ações que o Grupo tenha anunciado.

\begin{analysis}
\begin{itemize}[leftmargin=*,nosep]
\item \textbf{ODS 13: focar os fornecedores.} Mais de 95\% das emissões de gases com efeito de estufa do Grupo estão na cadeia de valor a montante e a jusante, e o Grupo já envolve os fornecedores em metas com base científica \cite[p.~42]{AR}. Reforçar a descarbonização dos fornecedores atuaria onde está a maior parte das emissões.
\item \textbf{ODS 3: fechar a lacuna de equidade.} A ambição é alcançar 400 milhões de pessoas de comunidades carenciadas até 2030, e 156 milhões tinham sido alcançadas desde 2024 \cite[pp.~5, 41]{AR}. Concentrar os programas de acesso nos grupos carenciados aceleraria o progresso rumo a essa meta.
\item \textbf{ODS 6, 12 e 15: impacto ambiental dos medicamentos.} O AR afirma que os resíduos farmacêuticos que entram no ambiente são atualmente um resultado inevitável do uso de medicamentos pelos doentes, e reporta medidas como avaliações de risco ambiental \cite[p.~45]{AR}. Estender essas medidas ao longo da cadeia de abastecimento é um próximo passo possível.
\end{itemize}
\end{analysis}
